% =============================================================================
% Ringkasan dan Summary -- Tesis Magister
%
% Berbeda dengan abstrak program sarjana, ringkasan tesis disusun dalam
% beberapa paragraf, dengan spasi tunggal, dan panjangnya tidak lebih dari dua
% halaman. Ringkasan memuat latar belakang permasalahan, tujuan, metode, hasil
% dengan penekanan pada temuan baru, dan implikasi. Tidak diperbolehkan
% mengacu pustaka, gambar, dan tabel.
%
% Ganti teks panduan pada ringkasan-id.tex dan ringkasan-en.tex dengan narasi
% ringkasan yang sebenarnya. Kata kunci diatur di src/config/information.tex.
% =============================================================================

\ipbprelimheading{RINGKASAN}
% Ganti teks panduan ini dengan narasi ringkasan dalam bahasa Indonesia.
% Ringkasan disusun dalam beberapa paragraf, spasi tunggal, maksimum dua
% halaman, dan tidak mengacu pustaka, gambar, atau tabel.
\noindent
\MakeUppercase{\NamaPenulis}. \JudulTesis\unskip. Dibimbing oleh
\DaftarPembimbingKapital

Narasi ditulis dalam satu spasi, disusun dalam beberapa paragraf, dan tidak
lebih dari dua halaman. Ringkasan/\textit{summary} memuat latar belakang
permasalahan, tujuan, metode, hasil dengan penekanan pada temuan baru, dan
implikasi yang disajikan secara ringkas, informatif, dan faktual. Tidak
diperbolehkan mengacu pustaka, gambar, dan tabel. Singkatan hanya dikenalkan
jika masih digunakan lagi dalam bagian lain dari ringkasan/\textit{summary}.
Ringkasan dalam bahasa Inggris ditulis dengan huruf miring
(\textit{italic}).


\ipbprelimsubheading{\itshape SUMMARY}
% Replace this guidance text with the English summary narrative.
% The summary is written in several paragraphs, single-spaced, at most two
% pages, and must not refer to literature, figures, or tables.
\begin{otherlanguage}{english}
\itshape
\noindent
\MakeUppercase{\NamaPenulisInggris}. \JudulTesisInggris\unskip. Supervised by
\DaftarPembimbingKapitalInggris

The narrative is written in single spacing, arranged in several paragraphs,
and not more than two pages. The summary contains the background of the
problem, objectives, methods, results with emphasis on new findings, and
implications presented concisely, informatively, and factually. It must not
refer to literature, figures, or tables. Abbreviations are introduced only if
they are used again in another part of the summary. The English summary is
written in italics.
\end{otherlanguage}


\vspace{0.48cm}

\noindent\hangindent=2cm\hangafter=1
Kata kunci: \KataKunci

\begin{otherlanguage}{english}
\itshape
\noindent\hangindent=2cm\hangafter=1
Keywords: \Keywords
\par
\end{otherlanguage}
